\begin{abstract}
Centralized school choice systems are increasingly adopted worldwide, yet their longer-term consequences for students' educational trajectories remain poorly understood. This paper examines the causal effect of implementing a centralized school assignment system in secondary education on subsequent college admissions outcomes. We exploit the staggered, region-by-region rollout of Chile's centralized school choice system, which provides quasi-experimental variation in exposure across cohorts, and estimate dynamic treatment effects to accommodate variation in treatment timing. Linking administrative records from secondary and postsecondary education, we find that centralized school choice significantly increases the likelihood that students apply to, are admitted to, and enroll in programs within the centralized college admissions system. These gains are concentrated among female students, students from socioeconomically vulnerable schools, and students from lower-performing schools. We further show that the improvements are driven by admissions to less selective and historically under-subscribed programs, with some displacement in highly selective programs. The results are robust across specifications, placebo tests, and sample restrictions. Together, our findings demonstrate that centralized school choice can meaningfully expand access to higher education, particularly for students historically underserved by decentralized admissions.
\end{abstract}
